\documentclass[12pt, a4paper]{article}

\usepackage[utf8]{inputenc}
\usepackage[T1]{fontenc}
\usepackage[italian]{babel}
\usepackage{amsmath, amssymb}
\usepackage{geometry}
\usepackage{booktabs}
\usepackage{array}       % modificatori colonna >{\raggedright}
\usepackage{xcolor}
\usepackage{parskip}
\usepackage{hyperref}
\usepackage{xurl}        % spezza URL a /, ., _, - ecc. (caricato dopo hyperref)

\geometry{margin=2.5cm}
\setlength{\emergencystretch}{3em}   % flessibilità extra contro overfull hbox

% Tipo colonna comodo: paragrafo con testo a sinistra
\newcolumntype{L}[1]{>{\raggedright\arraybackslash}p{#1}}

\hypersetup{
    colorlinks=true,
    linkcolor=blue!70!black,
    urlcolor=blue!70!black,
}

\title{\textbf{Fonti dei dati}\\[0.3em]
       \large Progetto COVID-19 — MSN}
\author{}
\date{}

\begin{document}

\maketitle
\tableofcontents
\newpage

% ---------------------------------------------------------------------------
\section{\texttt{iss-covid19.csv} — Casi per data di prelievo/diagnosi (ISS)}

\textbf{Sito:}
\url{https://www.epicentro.iss.it/coronavirus/sars-cov-2-sorveglianza-dati}

\textbf{File originale:}
\url{https://www.epicentro.iss.it/coronavirus/open-data/covid_19-iss.xlsx}\\
(foglio \texttt{casi\_prelievo\_diagnosi}, colonne: \texttt{DATA\_PRELIEVO} e \texttt{CASI})

\textbf{Formato nel progetto:} CSV con separatore \texttt{;}, intestazione
\texttt{DATA\_PRELIEVO\_DIAGNOSI;CASI}, date nel formato \texttt{dd/MM/yyyy}.

\textbf{Copertura:} 29/01/2020 – 06/01/2026 (2171 righe). Il valore
\texttt{08/01/2026} presente nel file originale è la data di aggiornamento del
dataset, non un dato epidemiologico.

\textbf{Note:} I casi sono per \textbf{data di prelievo} (o data diagnosi).
Questa è più vicina alla data reale di infezione rispetto alla data di referto
DPC, ma più tardiva rispetto alla data di inizio sintomi usata da ISS per il
calcolo ufficiale di Rt.

\textbf{Fogli presenti nell'xlsx (verificato):}

% Larghezze: 3.8 + 5.0 + 5.0 = 13.8cm  (testo ~16cm, spazio colonne ~0.85cm)
\begin{table}[h!]
\small
\begin{tabular}{@{} L{3.8cm} L{5.0cm} L{5.0cm} @{}}
\toprule
\textbf{Foglio} & \textbf{Contenuto} & \textbf{Variabili chiave} \\
\midrule
\texttt{casi\_prelievo\_diagnosi}
    & Casi per data prelievo/diagnosi (\textbf{usato})
    & \texttt{DATA\_PRELIEVO\_DIAGNOSI}, \texttt{CASI} \\[4pt]
\texttt{casi\_regioni}
    & Casi cumulativi per regione/PA
    & \texttt{NOME\_NUTS2}, \texttt{CASI\_CUMULATIVI} \\[4pt]
\texttt{casi\_province}
    & Casi cumulativi per provincia di residenza
    & \texttt{COD\_PROV}, \texttt{PROV\_names},
      \texttt{CASI\_CUMULATIVI} \\[4pt]
\texttt{ricoveri}
    & Ricoveri per data di ricovero
    & \texttt{DATARICOVERO1}, \texttt{RICOVERI} \\[4pt]
\texttt{decessi}
    & Decessi per data di decesso
    & \texttt{DATA\_DECESSO}, \texttt{DECESSI} \\[4pt]
\texttt{sesso\_eta}
    & Casi/decessi per sesso e fascia d'età
    & \texttt{SESSO}, \texttt{AGE\_GROUP},
      \texttt{CASI\_CUMULATIVI}, \texttt{DECEDUTI} \\
\bottomrule
\end{tabular}
\end{table}

\textbf{Il foglio con i casi sintomatici per data di inizio sintomi NON è
presente.} ISS non pubblica quei dati come open data. Il proxy
\texttt{nuovi\_positivi} DPC è quindi l'unica alternativa ragionevole per il
Metodo 2 del progetto.

\textbf{Archivi storici per anno:}
\begin{itemize}
    \item 2025: \url{https://www.epicentro.iss.it/coronavirus/open-data/OPENDATA-2025.zip}
    \item 2024: \url{https://www.epicentro.iss.it/coronavirus/open-data/OPENDATA-2024.zip}
    \item 2023: \url{https://www.epicentro.iss.it/coronavirus/open-data/OPENDATA-2023.zip}
    \item 2022: \url{https://www.epicentro.iss.it/coronavirus/open-data/OPENDATA-2022.zip}
    \item 2021: \url{https://www.epicentro.iss.it/coronavirus/open-data/OPENDATA-2021.zip}
    \item 2020: \url{https://www.epicentro.iss.it/coronavirus/open-data/OPENDATA-2020.zip}
\end{itemize}

% ---------------------------------------------------------------------------
\section{\texttt{dpc-covid19.json} — Andamento nazionale (Protezione Civile)}

\textbf{Repository GitHub:}\\
\url{https://github.com/pcm-dpc/COVID-19/blob/master/dati-json/dpc-covid19-ita-andamento-nazionale.json}

\textbf{Formato:} array JSON, un oggetto per ogni giorno. Campi rilevanti per
il progetto:
\begin{itemize}
    \item \texttt{data}: timestamp nel formato \texttt{yyyy-MM-ddTHH:mm:ss}
          (va troncato al giorno)
    \item \texttt{nuovi\_positivi}: nuovi positivi per \textbf{data di
          referto} — usato come $C(t)$ nel Metodo 1 (DPC) e Metodo 2
    \item \texttt{totale\_positivi}: attualmente positivi — usato come $I(t)$
          diretto nel Metodo 1 (DPC)
\end{itemize}

\textbf{Copertura:} 24/02/2020 in poi.

\textbf{Note:} Nel dataset scaricato \texttt{nuovi\_positivi} non presenta
valori negativi; il campo con valori negativi è
\texttt{variazione\_totale\_positivi} (differenza giornaliera di
\texttt{totale\_positivi}), che il codice non legge. Il \texttt{max(np, 0)}
presente in \texttt{carica\_dati.m} è comunque una misura difensiva: nella
storia del repository DPC ci sono stati episodi in cui correzioni retroattive
hanno prodotto \texttt{nuovi\_positivi} negativi (es.\ storno di duplicati
regionali), e la clausola garantisce robustezza qualora il file venisse
aggiornato con versioni precedenti o versioni future del dataset.

% ---------------------------------------------------------------------------
\section{\texttt{infn-rt.csv} — Rt ufficiale INFN CovidStat}

\textbf{Sito:} \url{https://covid19.infn.it/sommario/rt.html}\\
(dati estratti dal grafico interattivo della pagina — il certificato TLS del
dominio \texttt{covid19.infn.it} risulta non valido da browser, ma la pagina è
accessibile manualmente)

\textbf{Pagina metodologia:} \url{https://covid19.infn.it/sommario/rt-info.html}

\textbf{Pagina dati ISS (grafici Rt da sintomatici):}
\url{https://covid19.infn.it/iss/}

\textbf{Formato:} CSV con separatore \texttt{;}, decimale \texttt{,}. Colonne:
\begin{center}
{\small\ttfamily
giorno; rt; stat95\_upper; stat95\_lower; ci68\_lower;\\[2pt]
ci68\_upper; ci95\_lower; ci95\_upper}
\end{center}

\textbf{Copertura:} 18/03/2020 – 30/03/2023.

\textbf{Note:} Il progetto usa le colonne \texttt{rt} (valore centrale),
\texttt{ci95\_lower} e \texttt{ci95\_upper} (intervallo di confidenza al 95\%).

% ---------------------------------------------------------------------------
\section{Metodologia INFN/ISS per il calcolo di Rt}

\textbf{Algoritmo:} EpiEstim (pacchetto R) — Cori et al.,
\textit{American Journal of Epidemiology}, 2013.\\
DOI: \url{https://doi.org/10.1093/aje/kwt133}

\textbf{Blog tecnico INFN:}
\url{https://covidblog.infn.it/rt-e-linferenza-bayesiana-questa-sconosciuta/}

\subsection{Formula base (identica al Metodo 2 del progetto)}

\begin{equation}
    \mathbb{E}[I(t)] = R(t) \sum_{s=1}^{S} w(s)\, I(t-s)
\end{equation}

dove $w(s)$ è la distribuzione del tempo di generazione (= $\varphi(s)$ nel
codice).

\subsection{Distribuzione dei tempi di generazione (Cereda et al.)}

Distribuzione Gamma con parametri: \textbf{shape} $\alpha = 1.87$
(SD $= 0.26$), \textbf{rate} $\beta = 0.28$ (SD $= 0.04$) $\Rightarrow$
scala $\theta = 1/\beta \approx 3.57$ giorni, media $= \alpha\theta \approx
6.6$ giorni.

Fonte: Cereda et al. — esattamente i parametri usati in \texttt{phi\_gamma.m}.

\subsection{Stima bayesiana: EpiEstim vs Metodo 2 del progetto}

% Larghezze: 3.5 + 5.5 + 5.0 = 14.0cm
\begin{table}[h!]
\small
\begin{tabular}{@{} L{3.5cm} L{5.5cm} L{5.0cm} @{}}
\toprule
\textbf{Aspetto}
    & \textbf{EpiEstim (INFN/ISS ufficiale)}
    & \textbf{Metodo 2 (progetto, Metropolis)} \\
\midrule
Finestra di stima
    & 7 giorni scorrevoli (sliding window)
    & 1 giorno (per singolo $t$) \\[4pt]
Posteriore
    & Gamma analitica (coniugata Poisson-Gamma)
    & Campionata con MCMC (Metropolis) \\[4pt]
Prior su $R(t)$
    & Gamma(5, 5) — configurabile
    & Nessuno (MLE) \\[4pt]
Implementazione
    & R — pacchetto EpiEstim
    & MATLAB/Octave — da zero \\
\bottomrule
\end{tabular}
\end{table}

\subsection{Posteriore analitica di Cori}

Con prior $\mathrm{Gamma}(a, b)$ e likelihood Poisson su una finestra
$[t-\tau+1,\, t]$, la posteriore è esattamente:

\begin{equation}
    R(t)\mid\text{dati}
    \;\sim\;
    \mathrm{Gamma}\!\left(
        a + \sum_{k=t-\tau+1}^{t} I(k),\quad
        b + \sum_{k=t-\tau+1}^{t} \Lambda(k)
    \right)
\end{equation}

dove $\Lambda(k) = \sum_s w(s)\, I(k-s)$. Non serve MCMC: è analitica.

Il Metodo 2 del progetto approssima la stessa posteriore con $\tau=1$
(singolo giorno) tramite Metropolis. Per grandi $I(t)$ le due stime
convergono.

\subsection{Dati usati da ISS per Rt ufficiale}

ISS usa esclusivamente i \textbf{casi sintomatici} per \textbf{data di inizio
sintomi}, non i nuovi positivi DPC (data referto) né i casi ISS per data di
prelievo. Questo riduce il ritardo tra infezione e notifica e produce la stima
più accurata temporalmente.

\end{document}
